% =====================================================================
%  Chapter 8 -- Discussion
% =====================================================================
\chapter{Discussion}\label{ch:discussion}
\chapterlists[t]

\section{The three criteria, revisited}\label{sec:criteria-revisited}

\Cref{tab:criteria-evidence} collects the evidence for the criteria of \cref{sec:criteria}.

\begin{table}[!htb]
  \centering
  \small
  \renewcommand{\arraystretch}{1.3}
  \begin{tabularx}{\linewidth}{@{}>{\raggedright\arraybackslash}p{3.2cm}>{\raggedright\arraybackslash}X>{\raggedright\arraybackslash}p{2.1cm}@{}}
    \toprule
    \textbf{Criterion} & \textbf{Evidence} & \textbf{Where} \\
    \midrule
    1. Lag behind a moving threshold &
      Switch-on above and switch-off below threshold, in simulation and on the bench; hysteresis loops
      widening with rate. &
      \cref{fig:s1-traces,fig:s2-thresholds,fig:s2-loops} \\
    2. Lag grows predictably with rate &
      Switch-off shift $\propto\sqrt{\lambda}$ in simulation (exponent $0.499\pm0.003$) and on the bench
      ($0.56\pm0.08$), $k_{\mathrm{off}}$ within a few per cent of \eqref{eq:Roff}, fixed delay disfavoured.
      Switch-on exponent slightly above $\tfrac12$. In configuration B the law fits only marginally better
      than a fixed delay, with a coefficient $27\%$ below theory. &
      \cref{tab:s3-compare}, \cref{fig:s1-thresholds,fig:s3-thresholds,fig:s4-switchoff} \\
    3. Long-time state set by rate alone &
      On one fixed path, 10 runs re-ignite every cycle and 9 stay off, separated by
      $\SI{6.06}{\kilo\ohm\per\second}<\lamc<\SI{6.20}{\kilo\ohm\per\second}$; reproduced in simulation;
      absent for a single sweep. &
      \cref{fig:s4-twofates,fig:s4-cyclemap,fig:s4-tipping,fig:s4-single} \\
    \bottomrule
  \end{tabularx}
  \caption[The three criteria and the evidence for each]{The three criteria and the evidence.}
  \label{tab:criteria-evidence}
\end{table}

\section{One idea behind every result}\label{sec:one-idea}

Near threshold the amplitude changes exponentially at a rate proportional to the distance from threshold
(\cref{eq:sigma}), and switching state needs a fixed accumulated growth. Everything follows from how much
time the sweep allows to collect it:
\begin{itemize}
  \item one crossing: growth $\propto\lambda t^{2}$, giving the $\sqrt{\lambda}$ lag
        (\cref{eq:Ron,eq:Roff});
  \item a finite range: the lag eventually exceeds it, giving the first critical rate (\cref{eq:s2-lam1});
  \item repeated crossings: budgets $\propto 1/\lambda$ (\cref{eq:budget-rate}); when the decay budget runs
        out the circuit stays on (\cref{eq:s3-lamc}), when the growth budget runs out it stays off
        (\cref{eq:s4-lamc}).
\end{itemize}
The path decides which of the last two happens; the rate, threshold and (for tipping to quiet) the noise
decide where. Two predictions, $k_{\mathrm{off}}$ of \cref{eq:Roff} and $A_{\min}=\Asat e^{-\Gm}$ of
\cref{eq:Amin}, contain no free parameter and are reproduced to within a few per cent
(\cref{fig:s1-thresholds,fig:s3-amin}).

\section{What kind of tipping is this?}\label{sec:what-kind}

In its strict sense, R-tipping concerns a stable state that exists all along the path of the parameter
and is lost only because the system cannot follow it~\cite{Ashwin2012,Wieczorek2023}. Our circuit differs:
each pass crosses the Hopf threshold, so the quiet state is unstable above it and the oscillation does not
exist below it. What we observe is the \emph{signature} of R-tipping. The circuit is driven periodically
along a fixed path and ends in one of two long-time states, an oscillation rebuilt every cycle or a quiet
state never left, and which one is decided by the rate alone, across a sharp critical rate. The
single-sweep control shows that one pass cannot make this choice; it is the accumulation of many passes,
each too short for recovery. Noise only sets the level from which regrowth starts; at fixed noise the rate
alone decides. Experiments differing only in rate end in different states, which is the defining property
of R-tipping, here realised under periodic forcing across a bifurcation.

\section{The Wien-bridge oscillator as an analogue system}\label{sec:analogue}

The circuit is a well-characterised, controllable analogue system for rate-induced phenomena. Distance
from threshold, time scale ($RC$), path and rate, and potentially noise, are all independent circuit
parameters. The linear growth law holds until saturation, so lags, critical rates and residual amplitudes
can be predicted quantitatively, some with no free parameter. Successive cycles agree to tens of ohms,
bench and simulation agree once their thresholds are aligned, and the tipping boundary is sharp to $2\%$.
At \SI{73}{\hertz} a sweep cycle takes seconds, so one run contains dozens of realisations of the same
experiment, and component-level simulation separates the behaviour of nominal components from that of
the real ones.
